\documentclass[11pt]{article}
\usepackage{series}
\usepackage{booktabs}
\usepackage{array}

\title{\textbf{A Covariant Effective Model for Coherence-Capacity Cosmology}\\
\large Equations, Degeneracies, and Observational Tests}
\author{Peter Nero}
\date{Version 3, July 2026}

\hypersetup{
  pdftitle={A Covariant Effective Model for Coherence-Capacity Cosmology},
  pdfauthor={Peter Nero}
}

\newcommand{\CC}{\mathcal C}
\newcommand{\mpl}{M_{\mathrm{Pl}}}

\begin{document}
\maketitle

\begin{abstract}
This paper replaces an earlier heuristic cosmology with a covariant
effective-model interface.  Coherence capacity in Modal Triplet Theory
(MTT) is an admissibility margin; by itself it does not determine a
spacetime metric, expansion, acceleration, a causal horizon, a
cosmological singularity, or an arrow of time.  To ask those questions we
declare a Jordan-frame scalar--tensor action in which a scalar
$\CC$ is proposed as the large-scale image of capacity data.  We derive
the metric equation, scalar equation, homogeneous Friedmann equations,
and the exact condition for accelerated expansion.  We also prove a
reconstruction-degeneracy result: freely chosen coupling functions can
reproduce a broad class of expansion histories, so agreement with
$H(z)$ is not a prediction until MTT selects those functions and their
initial data independently of the observations.  Capacity-zero level
sets are separated from Lorentzian particle, event, and apparent
horizons; positive capacity is shown not to imply past completeness; and
time orientation is kept distinct from irreversible entropy production.
The result is a conditional effective realization, not a derivation of
cosmology from the present MTT axioms.  It supplies the equations and
data products that a future same-source theorem must determine.
\end{abstract}

\section*{Version 3 Revision Note}

\begin{description}
\item[Supersedes.] Version 2, \emph{Cosmology as Global Coherence Capacity
Evolution}.
\item[Reason.] The earlier paper inferred expansion, acceleration,
horizons, nonsingularity, and a time arrow from capacity language without
a covariant metric equation or a selected irreversible dynamics.
\item[Resolution.] Version 3 declares and varies a covariant effective
action, derives its homogeneous equations and exact acceleration
criterion, proves the reconstruction degeneracy and the relevant
separation counterexamples, and supplies an observational test contract.
\item[Retained content.] The proposal that an MTT admissibility margin may
have a useful large-scale cosmological image is retained as a conditional
interpretation.
\item[Open boundary.] MTT has not yet selected the source map, coupling
functions, initial state, perturbation branch, or a held-out
observationally viable solution from one accepted upper source.
\end{description}

\section{Introduction and orientation}

In plain language, the motivating picture is that a cosmological solution may carry, in
addition to its metric and matter fields, a slowly varying record of how
far the projected description lies from loss of admissibility.  It is
tempting to call that record ``spare capacity'' and then to read
expansion, acceleration, or horizons directly from it.  That temptation
must be resisted.  A positive scalar says that a margin is positive.  It
does not say how the metric evolves.

This revision therefore reverses the order of inference.  We first state
a covariant effective action.  Its Euler--Lagrange equations determine
which cosmological histories are allowed.  Only afterward do we ask
whether the scalar in that action can be sourced by the independently
defined MTT coherence capacity.  This places the proposal in the same
mathematical class as familiar scalar--tensor models
\cite{brans_dicke_1961,horndeski_1974,bellini_sawicki_2014}, while leaving
the MTT source problem visible.

The paper has three aims:
\begin{enumerate}
\item give a mathematically complete covariant realization that can be
reduced to homogeneous cosmology;
\item identify exactly which attractive conclusions do \emph{not} follow
from capacity language alone; and
\item specify the observables and independent tests required before the
realization can be promoted from a model to an MTT prediction.
\end{enumerate}

\section{Claim tier and interface with MTT}

\subsection{Three logically separate layers}

The construction uses three layers that must not be collapsed.

\begin{center}
\begin{tabular}{>{\raggedright\arraybackslash}p{0.19\linewidth}
                >{\raggedright\arraybackslash}p{0.34\linewidth}
                >{\raggedright\arraybackslash}p{0.36\linewidth}}
\toprule
Layer & Supplied object & What is not supplied \\
\midrule
MTT admissibility
& A normalized coherence-capacity margin and, where available, a
constitutive transport law
& A Lorentzian metric equation, a cosmological scalar, or an equation of
state \\
Covariant realization
& An action for $g_{\mu\nu}$, $\CC$, and matter
& A theorem selecting its free functions from upper geometry \\
Observational comparison
& Distances, expansion rates, perturbations, and likelihoods
& Predictive force unless parameters and functions were fixed before
using the compared data \\
\bottomrule
\end{tabular}
\end{center}

The current result is at the second layer.  The capacity papers define
and transport an admissibility margin
\cite{nero_capacity_2026,nero_capacity_dynamics_2026}.  They do not prove
that the resulting quantity is a fundamental scalar field, a varying
Planck mass, or a stress-energy component.  Those identifications are
hypotheses of the present effective model.

\subsection{What ``capacity cosmology'' means here}

\begin{definition}[Capacity-cosmology realization]
A capacity-cosmology realization is a tuple
\[
 \mathfrak C_{\mathrm{cos}}
 =
 (M,g,\CC,F,K,U,S_{\mathrm m},\mathfrak s)
\]
where $(M,g)$ is a time-oriented Lorentzian spacetime, $\CC$ is a real
scalar, $F,K,U$ are coupling functions, $S_{\mathrm m}$ is a matter
action, and $\mathfrak s$ is a proposed source map from declared MTT
capacity data to $\CC$.  The realization is \emph{selected} only if
$\mathfrak s$, $F$, $K$, $U$, and the relevant initial state are derived
from one accepted MTT source without fitting the target cosmological
observables.
\end{definition}

This definition makes room for the physical idea without pretending that
the source map has already been proved.

\section{A covariant effective action}

We use metric signature $(-,+,+,+)$ and units $c=1$.  Matter is minimally
coupled to the Jordan-frame metric $g_{\mu\nu}$, which therefore defines
redshift, rods, clocks, and the observational distances in this paper.
Consider
\begin{equation}
 S_{\mathrm{eff}}
 =
 \int_M \dd^4x\,\sqrt{-g}\,
 \left[
 \frac{1}{2}F(\CC)R
 -\frac{1}{2}K(\CC)
    g^{\mu\nu}\nabla_\mu\CC\nabla_\nu\CC
 -U(\CC)
 \right]
 +S_{\mathrm m}[g,\Psi].
 \label{eq:action}
\end{equation}
We assume $F(\CC)>0$.  A sufficient simple stability choice is
$K(\CC)\geq 0$; more generally the Einstein-frame scalar kinetic
coefficient must satisfy
\[
 \frac{K}{F}+\frac{3}{2}\left(\frac{F'}{F}\right)^2>0.
\]
The functions $F$, $K$, and $U$ are part of the model data, not results
of covariance.

\begin{theorem}[Field equations of the declared realization]
\label{thm:field-equations}
The stationary points of \cref{eq:action} satisfy
\begin{align}
 F G_{\mu\nu}
 &=
 T^{\mathrm m}_{\mu\nu}
 +K\left(
     \nabla_\mu\CC\nabla_\nu\CC
     -\frac{1}{2}g_{\mu\nu}(\nabla\CC)^2
   \right)
 -Ug_{\mu\nu}
 +\nabla_\mu\nabla_\nu F
 -g_{\mu\nu}\Box F,
 \label{eq:metric-field}\\
 0
 &=
 K\Box\CC
 +\frac{1}{2}K'(\nabla\CC)^2
 +\frac{1}{2}F'R
 -U'.
 \label{eq:capacity-field}
\end{align}
\par\medskip\noindent
If the matter equations hold, minimal coupling gives
$\nabla^\mu T^{\mathrm m}_{\mu\nu}=0$.
\end{theorem}

\begin{proof}
Metric variation of $\sqrt{-g}F R/2$ gives
$F G_{\mu\nu}/2+(g_{\mu\nu}\Box F-\nabla_\mu\nabla_\nu F)/2$
before the conventional factor of two is restored.  Variation of the
kinetic and potential terms gives the scalar stress tensor in
\cref{eq:metric-field}.  Varying $\CC$, integrating its kinetic term by
parts, and collecting the $F'R$ and $U'$ terms gives
\cref{eq:capacity-field}.  Diffeomorphism invariance and the matter
equations yield the matter conservation law.
\end{proof}

\paragraph{Interpretation.}
$F(\CC)$ is an effective squared Planck mass, $K$ controls the kinetic
response, and $U$ is an effective potential.  Calling $\CC$ a capacity
variable does not remove these ingredients.  If $\CC$ is instead a
nonlocal functional of more fundamental fields, varying
\cref{eq:action} as though it were independent would be incorrect; the
source map must then be inserted before variation.  That is one of the
open upstream choices.

\section{Homogeneous reduction and the acceleration test}

Take a spatially flat FLRW metric
\[
 \dd s^2=-\dd t^2+a(t)^2\dd\mathbf x^2,
 \qquad H=\frac{\dot a}{a},
\]
a homogeneous scalar $\CC(t)$, and a perfect fluid with density
$\rho_{\mathrm m}$ and pressure $p_{\mathrm m}$.  The field equations
reduce to
\begin{align}
 3F H^2
 &=
 \rho_{\mathrm m}
 +\frac{1}{2}K\dot{\CC}^{\,2}
 +U
 -3H\dot F,
 \label{eq:friedmann-one}\\
 -2F\dot H
 &=
 \rho_{\mathrm m}+p_{\mathrm m}
 +K\dot{\CC}^{\,2}
 +\ddot F-H\dot F,
 \label{eq:friedmann-two}\\
 0
 &=
 K(\ddot{\CC}+3H\dot{\CC})
 +\frac{1}{2}K'\dot{\CC}^{\,2}
 +U'
 -3F'(\dot H+2H^2),
 \label{eq:scalar-flrw}\\
 0
 &=
 \dot\rho_{\mathrm m}
 +3H(\rho_{\mathrm m}+p_{\mathrm m}).
 \label{eq:matter-continuity}
\end{align}
\par\medskip\noindent
These are genuine Friedmann-type equations because they are reductions
of a covariant action.  Dimensional consistency and isotropy alone would
not determine them.

\begin{theorem}[Exact acceleration criterion]
\label{thm:acceleration}
For every solution of
\cref{eq:friedmann-one,eq:friedmann-two}, accelerated expansion occurs
at a time $t$ if and only if
\begin{equation}
 2U
 >
 \rho_{\mathrm m}+3p_{\mathrm m}
 +2K\dot{\CC}^{\,2}
 +3\ddot F+3H\dot F.
 \label{eq:acceleration-condition}
\end{equation}
\end{theorem}

\begin{proof}
Multiply $\ddot a/a=H^2+\dot H$ by $6F$, then substitute twice
\cref{eq:friedmann-one} and three times
\cref{eq:friedmann-two}.  The result is
\[
 6F\frac{\ddot a}{a}
 =
 -\rho_{\mathrm m}-3p_{\mathrm m}
 -2K\dot{\CC}^{\,2}
 +2U-3\ddot F-3H\dot F.
\]
Since $F>0$, positivity of $\ddot a$ is equivalent to
\cref{eq:acceleration-condition}.
\end{proof}

The theorem identifies the error in the earlier heuristic.  An outward
capacity current, increasing $\CC$, or decreasing $1/\CC$ does not by
itself determine the signs of $\dot F$, $\ddot F$, $U$, or $\dot H$.
Acceleration follows only after the action and a solution fix all terms
in \cref{eq:acceleration-condition}.

\begin{corollary}[Minimal-coupling limit]
If $F=\mpl^2$ is constant, then
\[
 \frac{\ddot a}{a}>0
 \quad\Longleftrightarrow\quad
 2U>\rho_{\mathrm m}+3p_{\mathrm m}
       +2K\dot{\CC}^{\,2}.
\]
Thus a canonical kinetic term resists acceleration, while a sufficiently
large positive potential can support it.
\end{corollary}

\begin{proposition}[Capacity positivity does not imply expansion]
\label{prop:no-expansion}
Without a declared dynamical relation between $\CC$ and $g$, the
condition $\CC>0$ implies neither $H>0$ nor $\ddot a>0$.
\end{proposition}

\begin{proof}
The same constant positive function $\CC=\CC_0$ can be placed on a static
FLRW metric, an expanding de Sitter metric, or a contracting FLRW metric
at the purely kinematic level.  Capacity positivity distinguishes none
of them.  Within \cref{eq:action}, their admissibility is decided by
$F,K,U$, matter, and initial data rather than by the sign of $\CC$ alone.
\end{proof}

\section{Why a fitted expansion history is not yet a prediction}

A flexible scalar--tensor action can reconstruct a chosen background
history.  This is useful for model building and dangerous for claims of
derivation.

\begin{theorem}[Background reconstruction degeneracy]
\label{thm:reconstruction}
Let $H(t)$, $\rho_{\mathrm m}(t)$, and $p_{\mathrm m}(t)$ be smooth and
satisfy \cref{eq:matter-continuity}.  Choose a monotone clock
$\CC(t)$ and a smooth positive function $F(t)$.  Define
\begin{align}
 \mathcal K(t)
 &:=
 -2F\dot H-(\rho_{\mathrm m}+p_{\mathrm m})
 -\ddot F+H\dot F,
 \label{eq:reconstructed-kinetic}\\
 \mathcal U(t)
 &:=
 3F H^2-\rho_{\mathrm m}
 -\frac{1}{2}\mathcal K(t)+3H\dot F.
 \label{eq:reconstructed-potential}
\end{align}
If $\mathcal K(t)\geq0$, then along the selected trajectory the choices
\[
 K(\CC(t))=\frac{\mathcal K(t)}{\dot{\CC}(t)^2},
 \qquad
 U(\CC(t))=\mathcal U(t)
\]
reproduce that background.  Where $\dot{\CC}\neq0$, the scalar equation
follows from the metric equations, matter conservation, and the Bianchi
identity.
\end{theorem}

\begin{proof}
Substitution of \cref{eq:reconstructed-kinetic} into
\cref{eq:friedmann-two} makes the second Friedmann equation an identity.
Substitution of both reconstructed quantities into
\cref{eq:friedmann-one} makes the first equation an identity.  The
contracted Bianchi identity relates the divergence of the metric
equation to the scalar equation.  With conserved matter and
$\dot{\CC}\neq0$, the remaining scalar factor must vanish.
\end{proof}

\paragraph{Consequence.}
An observed $H(z)$ can be replayed by many choices of $F,K,U$.  Such a
reconstruction is not evidence that MTT selected the expansion history.
A prediction requires $F,K,U$, the source map, and initial data to be
fixed independently and then compared with held-out observations.

\subsection{A concrete example: the nested standard cosmology}

\begin{proposition}[$\Lambda$CDM is a nested limit]
\label{prop:lcdm}
For $F=\mpl^2$, constant $\CC$, and
$U=\mpl^2\Lambda$, \cref{eq:action} reduces at the background level to
general relativity with a cosmological constant:
\[
 3\mpl^2H^2=\rho_{\mathrm m}+\mpl^2\Lambda,
 \qquad
 -2\mpl^2\dot H=\rho_{\mathrm m}+p_{\mathrm m}.
\]
\end{proposition}

This limit is an important consistency check, but it is not a
capacity-based explanation of $\Lambda$.  A nonconstant alternative
must outperform or distinguish itself from this baseline after accounting
for all added functions and parameters.

\section{Transport data are not a metric equation}

The MTT transport program may supply an effective current
$J^\mu_{\CC}$ and a balance law
\begin{equation}
 \nabla_\mu J^\mu_{\CC}=\sigma_{\CC}.
 \label{eq:capacity-balance}
\end{equation}
For a homogeneous current $J^\mu_{\CC}=q_{\CC}u^\mu$,
\cref{eq:capacity-balance} becomes
\[
 \dot q_{\CC}+3Hq_{\CC}=\sigma_{\CC}.
\]
This equation describes dilution, production, or loss of the declared
charge density \emph{on a given cosmological geometry}.  It does not
replace \cref{eq:friedmann-one,eq:friedmann-two}.  In particular, the
$3Hq_{\CC}$ term already contains the expansion one might otherwise try
to derive from the balance law.

A future source theorem must specify whether $q_{\CC}$ equals $\CC$, is a
function of $\CC$, or is an independent transported density.  It must
also show how its stress tensor or its modification of $F,K,U$ follows
from the same action.  Treating a bookkeeping current as stress-energy
without that map would violate the type boundary.

\section{Horizons and singularities}

\subsection{Causal horizons}

Cosmological horizons are properties of Lorentzian causal structure.  In
an FLRW spacetime a particle-horizon radius and an event-horizon radius,
when the relevant integrals exist, are
\begin{align}
 R_{\mathrm p}(t)
 &=a(t)\int_{t_{\mathrm i}}^t\frac{\dd t'}{a(t')},\\
 R_{\mathrm e}(t)
 &=a(t)\int_t^{t_{\mathrm f}}\frac{\dd t'}{a(t')}.
\end{align}
For nonzero spatial curvature the apparent-horizon areal radius is
\[
 R_{\mathrm A}=\frac{1}{\sqrt{H^2+k/a^2}}.
\]
These notions are not interchangeable; event horizons are global, while
apparent horizons depend on a foliation.

\begin{proposition}[A capacity-zero surface is not a causal horizon]
\label{prop:horizon-separation}
A regular level set $\CC^{-1}(0)$ is not, from that fact alone, a
particle, event, or apparent horizon.  Conversely, a cosmological horizon
need not be a zero of $\CC$.
\end{proposition}

\begin{proof}
On de Sitter spacetime, a constant positive scalar $\CC=\CC_0$ coexists
with a cosmological event horizon \cite{gibbons_hawking_1977}; hence
capacity zero is not necessary.  On Minkowski spacetime, the scalar
$\CC=t$ has the regular zero set $t=0$, which is a spacelike Cauchy
surface rather than a causal horizon; hence capacity zero is not
sufficient.
\end{proof}

Capacity level sets may still be useful diagnostics of model breakdown.
Their geometric relation to a causal horizon requires an additional
theorem involving $g$, null expansions, or causal accessibility.  Area
and entropy statements then require the separately typed hypotheses
given in the horizon paper \cite{nero_horizons_2026}.

\subsection{Past completeness}

\begin{proposition}[Positive capacity does not remove a singularity]
\label{prop:singularity-separation}
The condition $\CC>0$ does not imply bounded curvature or geodesic
completeness.
\end{proposition}

\begin{proof}
Take a constant positive scalar on the spatially flat dust solution
$a(t)\propto t^{2/3}$ for $t>0$.  Its Ricci scalar diverges as
$t^{-2}$ toward $t=0$, while $\CC$ remains positive.  The capacity
condition does not control the curvature.
\end{proof}

A nonsingular or past-complete model must instead prove the relevant
extension and completeness conditions for a solution of
\cref{thm:field-equations}.  Bounded capacity, bounded $H$, or the absence
of a capacity-zero surface is not by itself such a proof.  Standard
geodesic-incompleteness results also show why accelerated expansion does
not automatically provide a past completion
\cite{borde_guth_vilenkin_2003}.

\section{Time orientation and irreversible evolution}

The action \cref{eq:action} is time-reversal invariant under the usual
transformation of a homogeneous solution.  It therefore does not select
an arrow of time.  A time orientation is part of the Lorentzian
background, whereas an entropy arrow requires an irreversible state,
boundary condition, coarse graining, or open-system dynamics.

If an entropy current $s^\mu$ is supplied and satisfies
\[
 \nabla_\mu s^\mu\geq0,
\]
then its integral can orient a thermodynamic arrow on the domain where
the inequality holds.  Likewise, a declared irreversible capacity law
could orient a capacity arrow.  Neither inequality follows from the sign
of $\CC$ or from the covariant action alone.  The phrase ``capacity
exhaustion'' is therefore a physical input until its generator and
entropy-production theorem are specified.

\section{Observable map and comparison contract}

\subsection{Background observables}

Given a solution and the normalization $a(t_0)=1$, redshift is
$1+z=1/a$.  In the spatially flat case the comoving and luminosity
distances are
\begin{equation}
 \chi(z)=\int_0^z\frac{\dd z'}{H(z')},
 \qquad
 d_L(z)=(1+z)\chi(z).
 \label{eq:distance-map}
\end{equation}
Useful derived quantities include
\[
 q(z)=-\frac{\ddot a}{aH^2},
 \qquad
 w_{\mathrm{eff}}(z)
 =-1-\frac{2\dot H}{3H^2}.
\]
Equations \cref{eq:friedmann-one,eq:friedmann-two,eq:scalar-flrw} and
\cref{eq:distance-map} are sufficient to produce predictions for
background expansion, supernova distances, and baryon-acoustic-oscillation
distance combinations after the matter and radiation sectors are fixed.

The standard cosmological data set is already restrictive.  CMB
anisotropies constrain the early background and perturbations
\cite{planck_2018_parameters}; supernovae constrain relative luminosity
distances \cite{riess_1998,perlmutter_1999,brout_2022}; and BAO constrain
radial and transverse distance scales.  A capacity model must be compared
with the full covariance information rather than with a few plotted
central values.

\subsection{Perturbations are indispensable}

Background agreement is not enough.  Perturbing
\cref{eq:metric-field,eq:capacity-field} determines gravitational slip,
the effective clustering strength, scalar propagation, and the growth of
matter perturbations.  The effective Planck-mass running
\[
 \alpha_M=\frac{\dot F}{HF}
\]
is one useful diagnostic, but it is not a complete perturbation theory.
The same selected functions used for $H(z)$ must also be used for CMB,
lensing, structure growth, local-gravity constraints, and gravitational
wave propagation.  Scalar--tensor effective parameterizations provide a
standard comparison language \cite{bellini_sawicki_2014}.

\subsection{Minimum reproducible comparison packet}

A claim of observational viability requires one versioned packet
containing:
\begin{enumerate}
\item explicit $F(\CC)$, $K(\CC)$, $U(\CC)$ and source-map formulas;
\item all dimensional normalizations, initial data, priors, and the
Jordan-frame convention;
\item a background and perturbation solver with numerical tolerances;
\item named data releases and complete likelihood definitions, including
covariances, nuisance parameters, and train/hold-out separation;
\item posterior or certified parameter intervals and residuals;
\item comparison with $\Lambda$CDM using the true parameter and function
count; and
\item hashes linking the source, solver, inputs, and output tables.
\end{enumerate}

No such fit is claimed in this paper.  The nested
\cref{prop:lcdm} shows that the action has a standard baseline.  It does
not show that a nontrivial capacity branch fits the data or predicts a
departure.

\section{Parameter accounting and falsifiability}

At its present tier the realization contains three functions
$F,K,U$, a source map $\mathfrak s$, and initial data.  A monotone field
redefinition can remove one functional redundancy locally, but at least
two independent functional degrees of freedom generally remain.  This is
more freedom than the six-parameter base $\Lambda$CDM background and
primordial model used in standard fits
\cite{planck_2018_parameters}.  The comparison is therefore not won by
renaming functions as geometry.

The model becomes sharply falsifiable after upstream selection.  Given
fixed source data and no target fitting, any of the following is an exit:
\begin{itemize}
\item $F\leq0$ or a negative Einstein-frame kinetic coefficient;
\item failure of the background equations or loss of a regular solution;
\item disagreement with distance, CMB, growth, lensing, local-gravity, or
gravitational-wave constraints outside the declared uncertainty budget;
\item a purported capacity horizon that fails the relevant causal or
null-expansion criterion; or
\item an arrow-of-time claim without a positive entropy-production law.
\end{itemize}

The decisive MTT theorem would derive
\[
 \mathfrak s,\quad F,\quad K,\quad U,\quad
 \text{initial state}
\]
from the same upper geometry that supplies the accepted capacity
operator, while preserving covariance and the current fixed-point and
operator interfaces.  Until then, \cref{eq:action} is a candidate
realization and not an inevitable consequence of MTT.

\section{Relation to established cosmology}

The equations in this paper are not a new class of gravitational theory.
They are a scalar--tensor effective action, a well-studied extension of
general relativity \cite{brans_dicke_1961,horndeski_1974}.  What is
specific to MTT is the proposed interpretation and the stronger demand
that the functions be selected by an independently defined
admissibility geometry.

This distinction provides a useful benchmark:
\begin{itemize}
\item If $F$ and $\CC$ are constant and $U=\mpl^2\Lambda$, the model is
the standard $\Lambda$CDM background.
\item If $F$ varies, the model belongs to the constrained
scalar--tensor/modified-gravity landscape.
\item If an MTT theorem fixes $F,K,U$ before looking at cosmological
data, the model acquires genuinely new predictive content.
\item If the functions are reconstructed from the observed $H(z)$, the
result is a profile or fit, not a derivation.
\end{itemize}

Capacity language may still prove useful by organizing which branch is
admissible and when an effective description exits its domain.  It does
not exempt the model from the Einstein equations, causal definitions, or
observational likelihoods appropriate to its declared action.

\section{Conclusion}

The corrected result is narrower and more useful.  A coherence-capacity
interpretation can be embedded in a covariant cosmological model, and
that model has explicit field equations, an exact acceleration test,
causal-horizon criteria, and an observational map.  None of these
structures follows from capacity positivity alone.

The principal mathematical lesson is the reconstruction degeneracy:
with freely chosen $F,K,U$, many expansion histories can be replayed.
The scientific frontier is therefore not to invent another
capacity-shaped Friedmann equation.  It is to derive the source map and
coupling functions from accepted MTT geometry before using cosmological
measurements, then run the same branch through background, perturbation,
causal, and statistical tests.  That would turn the present conditional
realization into a predictive cosmology.

\bibliographystyle{plain}
\bibliography{main}

\end{document}
